\subsection{Activity patterns (H2, H3)}\label{sec:r-activity}
\Cref{tab:burst} compares the activity statistics of both networks with P1.

\begin{table}[t]
\centering\small
\begin{tabular*}{\columnwidth}{@{\extracolsep{\fill}}llrr@{}}
\toprule Network & Metric & Real & P1 mean $\pm$ SD \\
\midrule
PS & $B$ & 0.535 & $0.5184\pm0.0004$ \\
PS & $M$ & 0.020 & $0.0710\pm0.0016$ \\
PS & Repeated ties & 0.625 & $0.6568\pm0.0001$ \\
PS & Jaccard & 0.220 & $0.3668\pm0.0010$ \\
CM & $B$ & 0.581 & $0.2451\pm0.0012$ \\
CM & $M$ & 0.100 & $0.2288\pm0.0057$ \\
CM & Repeated ties & 0.495 & $0.4949\pm0.0000$ \\
CM & Jaccard & 0.096 & $0.0690\pm0.0009$ \\
\bottomrule\end{tabular*}
\caption{Burstiness $B$, memory $M$, repeated-tie fraction and adjacent-window Jaccard
similarity (1-h windows for PS, 7-day windows for CM). The CM repeated-tie fraction is
invariant under P1 and therefore descriptive.}\label{tab:burst}
\end{table}

\paragraph{Burstiness and memory (H2).}
Burstiness exceeds the P1 value in both networks, by $\Delta B=+0.016$ in PS and $+0.336$
in CM, whereas memory is lower than under P1 in both ($\Delta M=-0.051$ and $-0.129$). The
PS excess is small in absolute terms over the full interval, although large relative to
the narrow null dispersion; it increases to $+0.116$ and $+0.120$ in the two single-day
halves (real/null 0.404/0.288 and 0.395/0.275). The inactive night between the two school
days is a plausible contributor but cannot be isolated with these data. The mean tie
lifetime is 0.444 days against 0.690 days under P1 in PS, and 4.756 days against 30.028
days in CM.

\paragraph{Persistence and turnover (H3).}
Adjacent-window similarity is lower than under P1 in PS ($\Delta J=-0.147$), indicating
rapid turnover of the contact set. In CM the full-period similarity exceeds that of P1
($\Delta J=+0.027$), but the sign reverses within both halves (0.074 against 0.113 and 0.118
against 0.140). The CM contrast thus depends on the observation interval: P1 redistributes
timestamps over the selected interval, so its baseline varies with the duration, activity
level and window composition of that interval, and the present comparison does not identify
which of these factors is responsible. The repeated-tie fraction of CM is invariant under
P1 by construction and serves only as a description.

\paragraph{Verdict on H2 and H3.}
H2 is supported for burstiness, with an interval-dependent magnitude in PS and a large
effect in CM, and is not supported for memory, which is weaker than under P1 in both
networks. For H3, PS shows rapid turnover, whereas the CM persistence contrast depends on
the observation interval and does not support a general claim of more persistent ties.
